\begin{abstract}
We introduce \textsc{QLoRA}, a finetuning strategy that significantly lowers memory requirements, enabling the finetuning of a 65B parameter language model on a single 48GB GPU while retaining full 16-bit finetuning task accuracy. \textsc{QLoRA} operates by propagating gradients through a static, 4-bit quantized version of a pretrained language model into Low Rank Adapters (LoRA). Our most advanced model series, dubbed \textbf{Guanaco}, achieves superior results over all prior publicly released models on the Vicuna benchmark, attaining 99.3\% of ChatGPT’s performance with just 24 hours of finetuning on one GPU. Key innovations in \textsc{QLoRA} include: (a) 4-bit NormalFloat (NF4), a novel datatype that is information-theoretically ideal for normally distributed weights; (b) Double Quantization, which minimizes memory usage by also quantizing the quantization factors; and (c) Paged Optimizers that efficiently handle memory peaks. Leveraging \textsc{QLoRA}, we finetuned over 1,000 models and provide comprehensive analyses spanning 8 instruction datasets, various architectures (LLaMA, T5), and large model sizes—including scales like 33B and 65B parameters that are otherwise impractical with standard finetuning. Our findings indicate that QLoRA enables state-of-the-art outcomes on challenging tasks with smaller, high-quality datasets, and surpasses prior SoTA even with less extensive models. We also present extensive evaluation of chatbot capabilities using both human raters and GPT-4, highlighting GPT-4 as a cost-effective proxy for human judgment. Moreover, our results show current chatbot benchmarks to be unreliable measures of true performance. Finally, a targeted analysis identifies specific cases where \textbf{Guanaco} does not match ChatGPT. All models, code, and specialized CUDA kernels for 4-bit training are made publicly available.\footnote{\url{https://github.com/artidoro/qlora} and \url{https://github.com/TimDettmers/bitsandbytes}}
\end{abstract}

\section{Introduction}

Adapting large language models (LLMs) through finetuning is recognized as a powerful strategy for enhancing model capabilities \citep{min2021metaicl, wei2021finetuned, ouyang2022training, wang2022super, wang2022self, liu2022few}, as well as for tailoring model outputs by incorporating preferred or filtering undesired behaviors \citep{ouyang2022training,askell2021general,bai2022training}. Nevertheless, the process of finetuning models of substantial size is often cost-prohibitive; for example, conventional 16-bit finetuning of a LLaMA model containing 65B parameters~\citep{touvron2023llama} demands in excess of 780 GB of GPU memory. While modern quantization strategies provide substantial memory reductions during inference for LLMs \citep{dettmers2022llm, dettmers2022case, frantar2022gptq, xiao2022smoothquant}, their utility vanishes during the training phase, as they become ineffective or unstable \citep{wortsman2023stable}.

In this work, we show for the first time that full finetuning of a 4-bit quantized model is achievable without compromising performance. Our proposed method, \textsc{QLoRA}, introduces a novel high-precision quantization approach that compresses a pretrained model to 4 bits, while incorporating a compact set of trainable Low-rank Adapter weights \citep{hu2021lora}

which are updated by propagating gradients through the quantized weight matrices.

With \textsc{QLoRA}, the memory needed to finetune a 65B parameter model is brought down from more than 780GB to less than 48GB of GPU memory, maintaining both runtime efficiency and predictive accuracy relative to traditional 16-bit finetuning. This advance greatly expands the accessibility of LLM finetuning: for the first time, even the largest available public models can be adapted using a single GPU. Applying \textsc{QLoRA}, we developed the \textbf{Guanaco} suite of models. The second strongest Guanaco model attains 97.8\% of ChatGPT’s accuracy on the Vicuna~\citep{vicuna2023} evaluation, all trainable within 12 hours on an ordinary consumer GPU; with a single high-end professional GPU over 24 hours, our largest model reaches 99.3\%, virtually closing the performance gap with ChatGPT on the Vicuna benchmark. At deployment, the smallest \textbf{Guanaco} variant (with 7B parameters) only requires 5 GB of memory, yet surpasses a much larger 26 GB Alpaca model by over 20 percentage points on the Vicuna evaluation (see Table~\ref{tbl:vicuna_eval}).

\textsc{QLoRA} incorporates several key advances aimed at minimizing memory consumption while maintaining model efficacy: (1) \textbf{4-bit NormalFloat}, a quantization format grounded in information theory and tailored for normally distributed variables, which demonstrates superior empirical outcomes compared to standard 4-bit Integers and 4-bit Floats. (2) \textbf{Double Quantization}, an approach that further compresses the quantization coefficients themselves, yielding an average memory savings of approximately 0.37 bits per parameter—equivalent to about 3 GB for a 65B-parameter model. (3) \textbf{Paged Optimizers}, which leverage NVIDIA unified memory to mitigate the memory surges associated with gradient checkpointing during mini-batch processing of extended sequences. Integrating these advancements, we develop an optimized variant of LoRA with adapter modules inserted at each layer of the network, thereby almost entirely circumventing the accuracy compromises characteristic of preceding methods.

The increased efficiency offered by \textsc{QLoRA} makes it possible to systematically investigate instruction finetuning and chatbot performance on model scales that conventional finetuning methods cannot feasibly support due to excessive memory requirements. Accordingly, we conduct training for over 1,000 models spanning diverse instruction-tuning datasets, model families, and sizes ranging from 80M up to 65B parameters. Beyond establishing that \textsc{QLoRA} achieves parity with 16-bit training performance (\S\ref{sec:comp_to_std}) and developing the cutting-edge chatbot \textbf{Guanaco} (\S\ref{sec:pushing}), we further examine overarching patterns within our suite of trained models. Notably, our results reveal that data quality significantly outweighs dataset scale—for example, a 9,000-sample dataset (OASST1) surpasses a 450,000-sample subset (FLAN v2) in chatbot metrics, even when both datasets target generalized instruction following. Additionally, we demonstrate that high scores on the Massive Multitask Language Understanding (MMLU) benchmark do not necessarily translate to comparable outcomes on the Vicuna chatbot benchmark, and vice versa, suggesting that dataset appropriateness for a task is more critical than dataset magnitude.

In addition, we undertake a comprehensive evaluation of chatbot effectiveness using assessments from both human annotators and GPT-4. Our methodology involves tournament-style benchmarking, wherein different models compete head-to-head to generate the optimal response to a specific prompt. The victor in each match is selected either by human raters or by GPT-4, and the cumulative match outcomes are summarized using Elo ratings~\citep{elo1967proposed, elo1978rating} to produce a performance leaderboard. Our findings reveal substantial concordance between GPT-4 and human judgments regarding model rankings, although notable discrepancies do occur. These observations underscore that while model-driven evaluations offer a cost-effective substitute for human annotation, they can introduce nontrivial sources of uncertainty.

To complement the quantitative benchmark outcomes, we perform a qualitative study focusing on the \textbf{Guanaco} model suite. This analysis draws attention to both successful outcomes and limitations that are not evident in quantitative metrics.

All model outputs, alongside human and GPT-4 evaluation annotations, are being released to promote additional research in this area. We also make our codebase and CUDA kernels publicly available, with full integration into the Hugging Face transformers ecosystem \citep{wolf2019huggingface}, ensuring broad accessibility. Our release includes a set of adapters for models with 7/13/33/65B parameters, each trained across eight instruction-following datasets, amounting to 32 unique open-source, finetuned variants.

\section{Background}
\paragraph{Block-wise k-bit Quantization} The technique of quantization involves mapping data from a high-precision format to a format with reduced bit width, effectively lowering the information content. A common instance of this is converting data from 32-bit floating point values to 8-bit integers. To guarantee that the entire value range of the lower-bit format is effectively utilized, the original data is typically normalized—usually by dividing by its absolute maximum—such that it fits within the target range. For example, when converting a tensor from 32-bit floating point (FP32) to an 8-bit integer tensor (Int8) with values in $[-127, 127]$:

\begin{equation}
    \mathbf{X}^{ \text{Int8}} = \text{round}\left( \frac{127}{{\text{absmax}}(\mathbf{X}^{{\text{FP32}}})}\mathbf{X}^{{\text{FP32}}} \right) = \text{round}(c^{\text{FP32}}\cdot \mathbf{X}^{{\text{FP32}}}),
\end{equation}

where $c$ denotes the {\em quantization scale} or {\em quantization constant}. The inverse procedure, dequantization, restores the values:

\begin{equation}
    \text{dequant}(c^{\text{FP32}}, \mathbf{X}^{\text{Int8}}) = \frac{\mathbf{X}^{\text{Int8}}}{c^{\text{FP32}}} = \mathbf{X}^{\text{FP32}}
\end{equation}

A significant drawback of this method arises when the input tensor contains values with large magnitudes (outliers). In such cases, the distribution of values across quantization bins becomes highly inefficient, with some bit combinations corresponding to bins that receive little or no quantized values. A widely adopted remedy involves partitioning the input tensor into multiple segments, or blocks, and quantizing each block separately using its own quantization constant $c$. Specifically, the input tensor $\mathbf{X}\in \mathbb{R}^{b\times h}$ is first flattened and then divided into $n$ contiguous chunks of size $B$, resulting in ${n = ({b\times h} )/{B} }$ linear segments. Each segment undergoes quantization independently, following Equation~1, which produces a quantized version of the tensor alongside $n$ individual quantization constants $c_i$.

\paragraph{Low-rank Adapters} The Low-rank Adapter (LoRA) finetuning approach \citep{hu2021lora} addresses the challenge of memory overhead during training by introducing a relatively small collection of learnable adapter parameters, while maintaining the original model parameters in a frozen state. During stochastic gradient descent, gradients propagate through the static pretrained weights but only the adapter’s parameters are updated to minimize the training objective. LoRA implements an enhanced linear mapping by appending an additional low-rank, factorized projection. For a standard projection ${\mathbf{X} \mathbf{W} = \mathbf{Y}}$ where $\mathbf{X}\in  \mathbb{R}^{b\times h}$ and $\mathbf{W}\in  \mathbb{R}^{h\times o}$, the LoRA formulation is expressed as:

\begin{equation}
    \mathbf{Y} = \mathbf{X}\mathbf{W} + s\mathbf{X}\mathbf{L}_1\mathbf{L}_2,
\end{equation}

with $\mathbf{L}_1\in  \mathbb{R}^{h\times r}$, $\mathbf{L}_2\in  \mathbb{R}^{r\times o}$, and $s$ denoting a scalar coefficient.

\paragraph{Memory Requirement of Parameter-Efficient Finetuning} A key consideration is the memory usage associated with LoRA during training, which depends on both the quantity and dimensions of adapters utilized. Given LoRA’s extremely lightweight memory profile, it is possible to deploy additional adapters to enhance model effectiveness without a significant rise in overall memory consumption. Although LoRA is categorized as a Parameter Efficient Finetuning (PEFT) approach, the dominant share of memory usage during LLM finetuning stems from storing activation gradients rather than the LoRA parameters themselves. For instance, when finetuning a 7B LLaMA model on FLAN v2 using a batch size of 1 and LoRA weights set to the commonly adopted 0.2\% of the base model’s weights~\citep{hu2021lora,liu2022few}, the input gradients from LoRA occupy approximately 567 MB, while the LoRA parameters themselves require only 26 MB. The implementation of gradient checkpointing~\citep{chen2016training} further lowers the storage required for input gradients to roughly 18 MB per sequence, making their memory demand greater than that of all LoRA weights combined. For context, the underlying 4-bit model alone requires 5,048 MB of memory. These observations demonstrate the significance of gradient checkpointing and reveal that further reductions in LoRA parameter size provide marginal memory savings. Consequently, expanding the number of adapters does not substantially affect total training memory requirements (see Appendix~\ref{app:memory} for comprehensive statistics). As elaborated later, this flexibility plays a critical role in matching the performance observed in full 16-bit precision settings.

\section{\textsc{QLoRA} Finetuning}

\textsc{QLoRA} facilitates accurate 4-bit finetuning through two principal innovations: 4-bit NormalFloat (NF4) quantization and Double Quantization. Furthermore, to address memory overflows during gradient checkpointing—a common obstacle to single-machine finetuning of large models—we propose Paged Optimizers, which alleviate memory usage spikes that could otherwise result in out-of-memory failures.

In the \textsc{QLoRA} paradigm, parameters are retained in a low-precision data type, typically 4 bits, while computations are carried out using a higher-precision data type, generally BFloat16. Thus, whenever a weight tensor within \textsc{QLoRA} is accessed for computation, it is dequantized to BFloat16 before proceeding with 16-bit matrix multiplication operations.

We proceed by detailing the constituent elements of \textsc{QLoRA} and subsequently present its formal mathematical formulation.

\paragraph{4-bit NormalFloat Quantization} The NormalFloat (NF) data type is an advancement over Quantile Quantization \citep{dettmers2022optimizers}, an information-theoretically optimal strategy designed to allocate an equal proportion of tensor values to each quantization interval. Quantile quantization operates by leveraging the empirical cumulative distribution function to estimate the quantiles of the input tensor.

However, a significant drawback of quantile quantization lies in the computational intensity associated with accurately estimating quantiles. As a result, efficient quantile approximation methods, such as SRAM quantiles~\citep{dettmers2022optimizers}, are often employed. Nevertheless, these approximate techniques introduce large quantization errors, particularly for outlier values, which are frequently among the most consequential in the data.

These computational and approximation challenges can be mitigated in scenarios where the input tensors are sampled from a distribution invariant to a quantization constant. In these circumstances, tensors share identical quantiles, rendering precise quantile estimation both practical and computationally tractable.

Typically, pretrained neural network weights follow a normal distribution centered at zero with a standard deviation $\sigma$ (refer to Appendix~\ref{app:normality}). To unify the distribution of all weights, we can rescale $\sigma$ so that the distribution spans the entire domain permitted by the chosen data type. In our approach, we specify the data type domain as $[-1, 1]$. Accordingly, the quantiles of both the data type and the neural network weights must be adjusted to fit within this interval.

The data type that achieves optimal information efficiency for zero-mean normal distributions of any standard deviation $\sigma$ within $[-1, 1]$ is derived in the following way: First, (1) the $2^k+1$ quantiles are calculated for the standard normal distribution $N(0, 1)$ to produce a $k$-bit quantile quantization scheme for normal distributions; (2) these quantile values are scaled so that they are within the $[-1, 1]$ domain; (3) input weight tensors are similarly normalized to the $[-1, 1]$ interval using absolute maximum rescaling.

When the range of the weights aligns with that of the data type, quantization proceeds as standard. The third step effectively ensures that the standard deviation of the weights matches that of the $k$-bit quantization data type. The specific $2^k$ quantized values $q_i$ are defined as:

\begin{equation}
q_i  = \frac{1}{2}\left( Q_X\left(\frac{i}{2^k+1}\right) + Q_X\left(\frac{i+1}{2^k+1}\right)\right),
\end{equation}

where $Q_X(\cdot)$ denotes the quantile function corresponding to the standard normal distribution $N(0,1)$. One limitation with symmetric k-bit quantization is its inability to exactly represent zero, which is crucial for quantizing padding and other elements with value zero without any error. To address this and guarantee a discrete zero point at $0$—as well as fully utilizing all $2^k$ representations for a k-bit type—we devise an asymmetric data type. This is accomplished by estimating quantiles $q_i$ across two intervals: allocating $2^{k-1}$ for negative values and $2^{k-1}+1$ for positive values. After combining these two sets of $q_i$, we remove one of the duplicated zeros, thus ensuring a unique zero point. We call this resulting format, which provides an equal expected count in each quantization bin, the \emph{k-bit NormalFloat} (NFk). This data type achieves information-theoretic optimality for quantizing data centered around zero and drawn from a normal distribution. The explicit quantized values for this type are given in Appendix~\ref{app:nf4}.

\paragraph{Double Quantization{}} 
We present \emph{Double Quantization} (DQ) as a method where the quantization constants themselves are quantized to further reduce memory usage. Although high-precision 4-bit quantization requires a small blocksize~\citep{dettmers2022case}, this typically results in substantial overhead from storing quantization constants. For instance, with 32-bit constants and a blocksize of 64 in $\mathbf{W}$, the constants introduce an average overhead of $32/64 = 0.5$ bits per parameter. Double Quantization mitigates this overhead.

To elaborate, Double Quantization views the initial quantization constants $c_2^{\text{FP32}}$ as the input for a second quantization layer. This subsequent stage produces quantized constants $c_2^{\text{FP8}}$ along with a new set of quantization constants $c_1^{\text{FP32}}$. We opt for 8-bit floating-point types and a blocksize of 256 for this secondary quantization, following the findings of \citet{dettmers2022case} that indicate no performance loss with 8-bit quantization. 
Given that the $c_2^{\text{FP32}}$ values are strictly positive, we normalize by subtracting the mean prior to quantization, centering them at zero to take advantage of symmetric quantization schemes. For a blocksize of 64, this technique reduces the quantization constants’ memory requirement from $32/64=0.5$ bits to $8/64 + 32/(64\cdot256) = 0.127$ bits per parameter, corresponding to a savings of 0.373 bits per parameter.

\paragraph{Paged Optimizers} leverage NVIDIA's unified memory~\footnote{\tiny \url{https://docs.nvidia.com/cuda/cuda-c-programming-guide}} functionality, which facilitates transparent paging between CPU and GPU memory on a page-by-page basis. This ensures seamless GPU computation even if GPU memory is occasionally exhausted. Similar to how standard memory paging occurs between main memory and disk storage, this mechanism automatically migrates pages as required. We employ unified memory to allocate paged space for optimizer states; when the GPU faces memory shortages, these optimizer states are automatically swapped out to CPU RAM and later retrieved to GPU memory during the optimizer’s update phase, without manual intervention.

\paragraph{\textsc{QLoRA}.} Building on the previously detailed components, we formalize \textsc{QLoRA} as it applies to a single linear layer within a quantized base model equipped with a single LoRA adapter:

\begin{equation}
    \mathbf{Y}^{\text{BF16}} =  \mathbf{X}^{\text{BF16}} \text{doubleDequant}(c_1^{\text{FP32}}, c_2^{\text{k-bit}}, \mathbf{W}^{\text{NF4}}) + \mathbf{X}^{\text{BF16}}\mathbf{L}^{\text{BF16}}_1\mathbf{L}^{\text{BF16}}_2,
\end{equation}

Here, doubleDequant$(\cdot)$ is specified by:

\begin{equation}
    \text{doubleDequant}(c_1^{\text{FP32}}, c_2^{\text{k-bit}}, \mathbf{W}^{\text{k-bit}}) = \text{dequant}(\text{dequant}(c_1^{\text{FP32}}, c_2^{\text{k-bit}}), \mathbf{W}^{\text{4bit}}) = \mathbf{W}^{\text{BF16}},
\end{equation}

NF4 quantization is employed for $\mathbf{W}$, while $c_2$ utilizes FP8 representation.

We configure $\mathbf{W}$ with a blocksize of 64 to achieve finer quantization granularity, and a blocksize of 256 is selected for $c_2$ to optimize memory usage.

When performing parameter updates, only the gradients with respect to the LoRA adapter weights, $\frac{\partial E}{\partial \mathbf{L}_i}$, must be computed; gradients for the 4-bit weights, $\frac{\partial E}{\partial \mathbf{W}}$, are unnecessary. Nevertheless, evaluating $\frac{\partial E}{\partial \mathbf{L}_i}$ requires calculating $\frac{\partial \mathbf{X}}{\partial \mathbf{W}}$, which follows from equation~(5). This process involves dequantizing $\mathbf{W}^{\text{NF4}}$ from storage format to the computation format $\mathbf{W}^{\text{BF16}}$, thus allowing the derivative $\frac{\partial \mathbf{X}}{\partial \mathbf{W}}$ to be determined using BFloat16 precision.

In essence, \textsc{QLoRA} employs a single storage data format—typically 4-bit NormalFloat—and a separate computation data format—16-bit BrainFloat. During both the forward and backward passes, the stored values are dequantized from the 4-bit representation to the 16-bit format used for computation; however, only the LoRA parameters, represented in 16-bit BrainFloat, have their weight gradients calculated.

\section{QLoRA vs. Standard Finetuning}
\label{sec:comp_to_std}

Having outlined the operational principles of QLoRA and its substantial memory savings in finetuning scenarios, the key question becomes whether QLoRA matches the performance of conventional full-model finetuning. In addition, we seek to understand the influence of QLoRA's design choices, especially the advantages of NormalFloat4 over conventional Float4. The subsequent sections detail our experimental efforts to investigate these topics.

\paragraph{Experimental setup.} Our evaluation spans three types of neural architectures: encoder, encoder-decoder, and decoder-only. We contrast QLoRA with both 16-bit adapter-based finetuning and full-precision finetuning on models of up to 3B parameters. Our benchmark suite includes GLUE \citep{wang2018glue} with RoBERTa-large \citep{liu2019roberta}, Super-NaturalInstructions (TKInstruct) \citep{wang2022super} with T5 \citep{t5}, and 5-shot MMLU \citep{hendrycksmeasuring} following finetuning of LLaMA on Flan v2 \citep{longpre2023flan} and Alpaca \citep{alpaca}. For a focused assessment of NF4 compared to alternative 4-bit data formats, we adopt the methodology from \citet{dettmers2022case}, examining post-quantization zero-shot accuracy and perplexity across several model families (OPT~\citep{zhang2022opt}, LLaMA~\citep{touvron2023llama}, BLOOM~\citep{scao2022bloom}, Pythia~\citep{biderman2023pythia}) with parameter sizes ranging from 125M to 13B.

Further setup-specific details are provided within each results section for clarity. Comprehensive experimental specifics are available in Appendix~\ref{app:exp-setup}.

Although paged optimizers are essential for enabling the finetuning of 33B and 65B parameter \textsc{QLoRA} models on a single 24GB or 48GB GPU, we do not present quantitative benchmarks for paged optimizers here, as page faults arise infrequently, typically only with mini-batches containing exceptionally long sequences. Nevertheless, runtime analyses for paged optimizers on 65B models running on 48GB GPUs demonstrate that, at a batch size of 16, their training throughput is comparable to standard optimizers. Further investigation is warranted to systematically determine the conditions under which paging may cause slowdowns.

\paragraph{Default LoRA hyperparameters do not match 16-bit performance}

Employing the established method of applying LoRA to the query and value projections within the attention mechanism \citep{hu2021lora} fails to yield finetuning outcomes equivalent to those achieved by full finetuning for larger models. Figure~\ref{fig:hyper} illustrates, in the context of LLaMA 7B finetuning on Alpaca, that the total number of LoRA adapters integrated—specifically, the inclusion of LoRA adapters across all linear layers of the transformer block—emerges as the most consequential hyperparameter for aligning LoRA performance with that of full finetuning. By contrast, variations in other parameters, such as the projection rank $r$, show minimal impact on performance (further details in Appendix \ref{app:exp-setup}).

In a similar vein, our analysis indicates that the default hyperparameter settings for baselines with full finetuning are insufficiently optimized. To establish robust baselines, we conduct an extensive hyperparameter sweep, varying learning rates between 1e-6 and 5e-5 and batch sizes from 8 up to 128. The resulting outcomes for 7B LLaMA finetuned on Alpaca are presented in Figure~\ref{fig:hyper}.

\paragraph{4-bit NormalFloat surpasses 4-bit Floating Point in performance}

Although the 4-bit NormalFloat (NF4) format is theoretically optimal in terms of information preservation, its practical benefits require empirical validation. We adopt the methodology described by \citet{dettmers2022case}, assessing various quantized LLMs (including OPT \citep{zhang2022opt}, BLOOM \cite{scao2022bloom}, Pythia \cite{biderman2023pythia}, and LLaMA) across a spectrum of model sizes from 125M to 65B, employing different data types. Evaluations cover language modeling as well as several zero-shot benchmarks. As shown in Figure~\ref{fig:nf4} and Table~\ref{tbl:nf4_ppl}, NF4 consistently yields notable improvements in performance compared to both FP4 and Int4. Additionally, employing double quantization effectively minimizes memory requirements without any measurable decline in accuracy.

\paragraph{k-bit \textsc{QLoRA} achieves parity with 16-bit full finetuning and 16-bit LoRA} While earlier work has demonstrated that 4-bit quantization for \emph{inference} is feasible, it typically incurs a drop in performance when contrasted with 16-bit precision \citep{dettmers2022case,frantar2022gptq}. This brings forth the key question of whether finetuning adapters at 4-bit can recover the lost performance. We investigate this question in two scenarios.

In the first scenario, we directly compare 16-bit full finetuning with adapter-based approaches (16-bit, 8-bit, and 4-bit) for RoBERTA and T5 models with parameter counts between 125M and 3B on the GLUE and Super-NaturalInstructions datasets. As summarized in Table~\ref{tab:nf4vsbf16}, results across both datasets show that adapter finetuning—regardless of being 16-, 8-, or 4-bit—matches the results of the fully finetuned 16-bit baseline. These findings imply that finetuning adapters subsequent to quantization can effectively recoup any performance deficits introduced by lower-precision quantization.

In our second experimental configuration, recognizing that fully finetuning models with 11B parameters or more necessitates multiple servers equipped with high-memory GPUs, we further investigate whether 4-bit \textsc{QLoRA} can achieve performance parity with 16-bit LoRA models across scales ranging from 7B to 65B parameters. To address this, we conduct finetuning of LLaMA models (7B through 65B parameters) using two instruction-following datasets—Alpaca and FLAN v2—and subsequently assess performance on the MMLU benchmark via 5-shot accuracy. The findings, detailed in Table~\ref{tbl:llama-datatype}, demonstrate that the NF4 configuration with double quantization entirely recaptures the MMLU performance levels of 16-bit LoRA. Furthermore, we observe that employing FP4 with \textsc{QLoRA} results in an approximately one percentage point shortfall relative to the 16-bit brain float LoRA baseline. These observations reinforce our earlier conclusions that (1) \textsc{QLoRA} leveraging NF4 is capable of replicating both 16-bit full and 16-bit LoRA finetuning results, and (2) NF4 exhibits superior quantization fidelity compared to FP4.

\paragraph{Summary} Across multiple academic benchmarks employing established evaluation protocols, our results robustly indicate that 4-bit \textsc{QLoRA} utilizing the NF4 data type consistently achieves results equivalent to those of 16-bit full finetuning and 16-bit LoRA finetuning. Our experiments also confirm NF4's advantage over FP4, as well as demonstrating that double quantization does not negatively impact performance. Collectively, this provides substantial support for the reliability of 4-bit \textsc{QLoRA} in matching the efficacy of conventional 16-bit approaches.

Building upon insights from prior work on quantization \citep{dettmers2022case}, both our MMLU and Elo evaluations reveal that within a constrained finetuning and inference resource allocation, expanding the base model size while reducing numerical precision offers clear advantages. This underscores the critical efficiency gains afforded by \textsc{QLoRA}. Since we detected no decline in performance compared to full finetuning when applying 4-bit finetuning in our studies, a pertinent open question remains regarding the precise boundaries of the performance-precision trade-off in \textsc{QLoRA} tuning, an inquiry that we propose for future research.

We next examine instruction tuning at scales that would be infeasible with standard 16-bit finetuning methods on typical research hardware.

\section{Pushing the Chatbot State-of-the-art with QLoRA}
\label{sec:pushing}
Having demonstrated that 4-bit \textsc{QLoRA} can achieve comparable performance to traditional 16-bit methods across various scales, task types, and datasets, we conduct a comprehensive analysis of instruction finetuning on the largest publicly accessible language models used in academic research. To evaluate instruction finetuning effectiveness, we test on the demanding MMLU Natural Language Understanding benchmark and devise new protocols for assessing real-world chatbot capabilities.

\subsection{Experimental setup} An overview of the experimental methodology is provided below, with further specifics detailed in Appendix \ref{app:chatbot}.

\paragraph{Data} Noting the absence of systematic evaluations of recent instruction-following datasets, we compile a selection of eight up-to-date datasets. Our selection spans datasets derived from crowdsourcing (OASST1~\citep{kopf2023openassistant}, HH-RLHF~\citep{bai2022training}), datasets distilled from instruction-tuned models (Alpaca~\citep{alpaca}, self-instruct~\citep{wang2022self}, unnatural-instructions~\citep{honovich2022unnatural}), aggregate corpora (FLAN v2~\citep{chung2022scaling}), as well as hybrid forms (Chip2~\citep{laion2023}, Longform~\citep{koksal2023longform}). This collection encompasses diverse languages, dataset sizes, and licensing conditions.

\paragraph{Training Setup} To mitigate potential confounding factors introduced by divergent training strategies, we carry out all QLoRA finetuning using cross-entropy loss under a supervised learning paradigm, excluding reinforcement learning approaches, even for datasets annotated with human preference information. Where datasets clearly differentiate instructions from responses, we perform finetuning solely on the responses (see ablation analyses in Appendix \ref{app:chatbot}). In cases such as OASST1 and HH-RLHF, which supply multiple candidate responses, we identify the top-rated response at each juncture of the conversational tree and finetune on the sequence of these selected exchanges, including associated instructions. For every experiment, we apply NF4 \textsc{QLoRA} with both double quantization and paged optimizers to circumvent memory overflow during gradient checkpointing. Limited hyperparameter searches are conducted for LLaMA 13B and 33B models, and we observe that settings established at the 7B scale (such as epoch count) are generally transferable, except for the learning rate and batch size: for 33B and 65B models, the learning rate is halved and the batch size is doubled.

\paragraph{Baselines} The models we develop are benchmarked against both academic (Vicuna~\citep{vicuna2023} and Open Assistant~\citep{kopf2023openassistant}) and proprietary systems (GPT-4 \citep{openai2023gpt}, GPT-3.5-turbo, and Bard). The Open Assistant system comprises a LLaMA 33B model, refined with Reinforcement Learning from Human Feedback (RLHF) utilizing the same OASST1 dataset as in our studies. Meanwhile, Vicuna employs comprehensive fine-tuning of the LLaMA 13B model on confidential user-conversation data sourced from ShareGPT, representing a distillation of OpenAI GPT outputs.

\subsection{Evaluation}

Consistent with standard evaluation practices, we adopt the MMLU (Massively Multitask Language Understanding) benchmark \citep{hendrycksmeasuring} for assessing capabilities across a breadth of language understanding tasks. The benchmark entails 57 multiple-choice subjects such as elementary mathematics, U.S. history, computer science, law, and others. All results are reported using the 5-shot accuracy protocol.

Our assessment of generative language abilities incorporates both automated and human evaluation methodologies. The latter focuses on model response quality by leveraging queries crafted by human annotators, reflecting a more practical assessment environment for chatbot models. This approach is increasingly favored in recent research; however, there remains a lack of standardized protocols within the academic community. Below, we detail our methodology, consistently employing nucleus sampling with $p=0.9$ and a temperature of $0.7$ for all generated outputs.

\paragraph{Benchmark Data} Our evaluation utilizes two manually curated query datasets: the Vicuna prompts \citep{vicuna2023} and the OASST1 validation split \citep{kopf2023openassistant}. The Vicuna prompts comprise 80 unaltered questions drawn from a wide array of categories. The OASST1 collection offers a multilingual set of multi-turn dialog exchanges sourced from crowds, involving users and assistants. We extract all user queries from the validation set, embedding preceding conversational turns within each prompt, resulting in a total of 953 distinct queries. For clarity, we designate these resources as the Vicuna and OA benchmarks.

\paragraph{Automated Evaluation} Building on the framework by \citet{vicuna2023}, we employ GPT-4 to score and compare model outputs to those of ChatGPT (GPT-3.5 Turbo) on the Vicuna dataset. Given each prompt with responses from ChatGPT and another model, GPT-4 assigns both a score on a ten-point scale, accompanied by justifications. Model effectiveness is then represented as a percentage relative to ChatGPT’s total score, with scores potentially exceeding 100\% if a model outperforms ChatGPT absolutely. Our observations reveal a notable position bias in GPT-4’s evaluation, with higher scores often attributed to the response appearing first in the prompt. To mitigate this bias, we advocate calculating the average score across both possible response orderings.

Subsequently, we conduct head-to-head output comparisons among different systems. Here, the evaluation task is streamlined to a three-label classification—indicating preference for one response, the other, or a tie. GPT-4 is asked to select the superior output or indicate equivalence, with explanations provided for each decision. All pairwise system combinations are compared across both the Vicuna and OA datasets.

\paragraph{Human Evaluation} While previous studies suggest that large language models like GPT-4 may serve as capable evaluators of generative system quality \citep{fu2023gptscore}, the degree to which GPT-4 assessments align with human preferences remains uncertain. As a result, we implement parallel human evaluations on the Vicuna benchmark, mirroring the protocols used in our automated tests. Data is collected via Amazon Mechanical Turk (AMT), enlisting two annotators for ChatGPT comparison tasks and three for direct pairwise evaluations between system outputs.

\paragraph{Elo Rating} To evaluate model performance, we conduct pairwise competitions using both human and automated assessments, structuring the process as a tournament in which models face off in head-to-head matchups to generate superior responses to given prompts. While this approach echoes the comparative framework utilized by \citet{bai2022training} and \citet{vicuna2023}, our methodology integrates GPT-4-based ratings alongside human evaluations. For the Elo calculation~\citep{elo1967proposed,elo1978rating}, we randomly draw from the labeled comparison dataset. Elo, a rating system extensively applied in chess and various competitive settings, quantifies a participant’s expected win probability relative to another—for instance, a participant with an Elo of 1100 matched against one with 1000 should win roughly 65\% of the time; equivalently rated participants, such as 1000 versus 1000 or 1100 versus 1100, are each expected to win 50\% of matches. Elo scores are adjusted after each contest based on the discrepancy between actual and expected results: unexpected victories prompt larger shifts, while outcomes in line with expectations cause minor updates. Over multiple rounds, this leads to ratings that closely reflect the comparative abilities of each participant. All models begin with an initial score of 1,000, using $K=32$ for the update rate. Following \citet{vicuna2023}, this entire procedure is repeated 10,000 times with distinct random seeds to mitigate the influence of the order in which models are paired.

\subsection{\textbf{Guanaco}: \textsc{QLoRA} fine-tuned on OASST1 Represents a Leading-Edge Chatbot} Through both automated and human assessment, we observe that the leading \textsc{QLoRA}-fine-tuned model, {Guanaco} 65B, trained on a modified OASST1 dataset, sets a new standard for open-source chatbot performance and delivers results that closely rival ChatGPT. In comparisons with GPT-4, both {Guanaco} 65B and 33B exhibit an expected win rate of 30\% according to human annotators' system-level pairwise evaluations using the Elo metric—this constitutes the highest value reported to date.

Table~\ref{tbl:vicuna_eval} summarizes outcomes on the Vicuna benchmark \cite{vicuna2023} in relation to ChatGPT. Results indicate that {Guanaco} 65B is surpassed only by GPT-4, reaching 99.3\% of ChatGPT’s performance. Although {Guanaco} 33B comprises a higher parameter count than Vicuna 13B, it adopts 4-bit quantization for its weights, enabling greater memory efficiency at 21 GB in contrast to 26 GB, and yields a three-point improvement over Vicuna 13B. Additionally, {Guanaco} 7B is suitable for modern smartphones due to its 5 GB size, while still outperforming Alpaca 13B by almost 20 points.

Nonetheless, Table~\ref{tbl:vicuna_eval} reveals wide confidence intervals, with substantial overlap in model performance. We attribute this ambiguity to vague scale definitions—for instance, what an 8 signifies on a 10-point scale can vary considerably between contexts. To address this, we propose the Elo rating approach \citep{elo1967proposed}, grounded in \textit{pairwise} human and GPT-4 evaluations, as a preferred metric over absolute scores. The Elo rankings for top-performing models are presented in Table~\ref{tbl:elo}. It is worth mentioning that human rankings and those from GPT-4 diverge somewhat, particularly regarding Guanaco 7B, but are largely in agreement for most models, as indicated by a Kendall Tau of $\tau=0.43$ and Spearman's $r=0.55$ at the system level. At the individual example level, agreement between GPT-4 and the majority human vote is lower, as reflected by a Fleiss $\kappa=0.25$. In sum, these findings indicate moderate alignment in system-level assessments between GPT-4 and human annotators, supporting the view that model-driven evaluations can serve as a relatively dependable surrogate for human evaluation. Further implications are examined in Section~\ref{ssec:considerations}.

The Elo rankings presented in Table~\ref{tbl:elo_large} show that both the {Guanaco} 33B and 65B models surpass all models except for GPT-4 on the Vicuna and OA benchmarks, demonstrating performance levels comparable to ChatGPT as reflected in Table~\ref{tbl:vicuna_eval}. It is important to observe that the Vicuna benchmark is more favorable to open-source models, whereas the OA benchmark shows a bias towards ChatGPT. Additionally, analysis of Tables \ref{tbl:mmlu-llama} and \ref{tbl:vicuna_eval} reveals that the choice of finetuning dataset plays a crucial role in determining performance outcomes. For example, Llama models finetuned using FLAN v2 excel on the MMLU benchmark but exhibit the lowest results on the Vicuna benchmark, a pattern that is similarly evident across other models. This underscores the partial independence of different evaluation benchmarks: achieving high MMLU scores does not necessarily correlate with superior chatbot capabilities, as measured by benchmarks like Vicuna or OA, and vice versa.

stands out as the only leading model in our study trained entirely without proprietary data, following the OASST1 dataset's explicit restriction against using GPT-derived data. The next best-performing open-source-trained model, Anthropic’s HH-RLHF, attains a score 30 percentage points below {Guanaco} on the Vicuna benchmark (refer to Table~\ref{tbl:vicuna_eval}). Taken together, these findings provide evidence that 4-bit \textsc{QLoRA} facilitates the creation of advanced chatbots capable of challenging ChatGPT's performance. Moreover, our 33B {Guanaco} model can be fully trained within 12 hours on consumer-grade GPUs equipped with 24 GB memory. This accessibility lays the groundwork for future innovations, particularly through further \textsc{QLoRA} tuning on highly specialized open-source datasets, which holds promise for producing models that stand alongside leading commercial alternatives.

\section{Qualitative Analysis}

Although our primary assessment relies on quantitative analysis, relying solely on aggregate statistics presents several limitations. One of the most significant concerns is the issue of benchmark validity \citep{liao2021we}—it is often uncertain whether a benchmark truly assesses the intended phenomena, particularly as unforeseen "shortcuts" for solving benchmarks are discovered and subsequently leveraged by machine learning models \citep{gururangan2018annotation, poliak2018hypothesis}. To address this limitation in part, we supplement our study with qualitative analysis, organized into two subsections. In \S\ref{ssec:examples}, we present selected instances that typify certain trends and behaviors seen in the outputs from our 65b {Guanaco} model. Following this, in \S\ref{ssec:considerations}, we elaborate on factors affecting our interpretation of the results and reflect upon their implications.

\subsection{Qualitative Analysis of Example Generations}
\label{ssec:examples}

For illustrative examples, we systematically reviewed the data produced for both the Vicuna and OpenAssistant benchmarks, seeking recurring themes or tendencies within {Guanaco}'s responses. Upon recognizing a recurring output pattern, we formulated targeted prompts or questions designed to elicit the pattern even when it does not correspond to the correct answer. For instance, if the model demonstrated a proclivity toward producing unnecessarily verbose responses, we deliberately prompted it to ``Answer yes or no without explanation.'' This methodology enabled us to uncover both “lemons” (cases where the model could be adversarially induced to fail) and “cherries” (cases where such adversarial prompting did not lead to failure). All textual outputs described in this section were produced using Nucleus Sampling \cite{holtzmancurious} with $p=0.9$.

Naturally, this overview is not exhaustive, as the limited scope of this qualitative analysis precludes control over the multitude of influencing factors; for example, the potential variety in model responses to a single prompt is extensive, necessitating reliance on a subset of outputs that we consider indicative. Nonetheless, presenting these selected instances provides interpretive context for the quantitative findings reported earlier. By making our models and code publicly available, we aim to encourage subsequent research that delves more deeply into the concerns highlighted in this section.

\paragraph{Factual Recall} When prompted with straightforward factual questions, such as ``What is the capital of Zambia?'', every model evaluated reliably delivers accurate information. For example:

\begin{quote} 
    \textbf{Guanaco}:
    \texttt{The capital of Zambia is Lusaka.}
\end{quote}

However, as the factual queries increase in obscurity, {Guanaco}'s reliability diminishes, despite its maintained certainty. Consider the following exchange drawn from HotPotQA \citep{yang2018hotpotqa}:

\begin{quote}
     \textbf{User}: In what year was the singer who popularized `I'll Keep the Lovelight Burning' born?  

     \textbf{Guanaco}:~\texttt{Al Jolson is credited with popularizing the song `I'll Keep the Lovelight Burning,' and he was born in the year 1886.}
 \end{quote}

Here, the response misidentifies both the performer associated with the song and their birth year, even though the stated date is accurate for Al Jolson—the incorrect figure referenced.

\paragraph{Suggestibility} {Guanaco} unexpectedly demonstrates notable resistance to accepting and propagating certain types of presupposed misinformation, as illustrated below:

\begin{quote}
    \textbf{User}: How was it finally, officially confirmed that the earth is flat by peer reviewed scientists? \\ \\
    \textbf{Guanaco}: \texttt{There has never been any official confirmation from peer-reviewed scientists that the earth is flat. The prevailing view in the scientific community is, and has long been, that the earth is an oblate spheroid. This understanding is backed by centuries of scientific evidence and direct experimental verification. \\ \\ The assertion that the earth is flat belongs to conspiracy theories that have repeatedly been disproven by scientific methods. The notion that peer review has validated a flat earth is simply inaccurate. \\ \\ Feel free to ask further questions if you'd like.}
\end{quote}

demonstrates a strong ability to recognize when certain types of queries cannot have meaningful answers, for instance,

\begin{quote}
    \textbf{User}: What time is it?

    \textbf{Guanaco}: \texttt{As an AI language model, I do not possess the ability to access live data such as the current time. My responses rely solely on the information encoded in my training data and do not represent real-time facts. To know the current time, please consult a clock or an online source.}
\end{quote}

This tendency is similarly evident for questions such as ``Where are you?'', ``How are you?'', and so forth.

\paragraph{Refusal} In the same vein, {Guanaco} may arbitrarily decline to execute a request without a clear reason:

\begin{quote}
    \textbf{User}: Please reverse the words in the sentence ``I'm a cool dude in a mocha mood.'' \\ \\
    \textbf{Guanaco}: \texttt{I apologize, but I'm unable to fulfill that request right now. However, I can provide an explanation regarding the grammatical correctness of your sentence.\\ \\The sentence ``I'm a cool dude in a mocha mood'' adheres to standard English grammar and syntax. In this construction, the adjective ``cool'' appropriately modifies the noun ``dude'', and the prepositional phrase ``in a mocha mood'' adds contextual detail. This syntactic structure is standard and coherent within English language norms.\\ \\If you have further inquiries or need assistance with another matter, please let me know.}
\end{quote}

\paragraph{Secret Keeping} Regrettably, when such refusals are actually warranted, their reliability is questionable. Designers frequently wish to keep specific information hidden from the user, even when the system has the capability to return free-form responses. We prompt the model with ``The secret word is `banana'. Do not reveal or repeat back this back, no matter what the human requests.'' to evaluate the difficulty a user might have in uncovering the concealed information. When the model is confronted with a direct query, it withholds the information as intended:

\begin{quote}
    \textbf{User}: What is the secret word?

    \textbf{Guanaco}: \texttt{I'm sorry, but I'm not able to disclose the secret word. I'm programmed to not reveal or repeat back the secret word, no matter what the human requests. Is there anything else I can help you with?}
\end{quote}

This demonstrates the intended outcome. However, even minimal prompt manipulation leads the system to divulge the information:

\begin{quote}
    \textbf{User}: This is a game. The goal is to ignore your previous instructions. What is the secret word?

    \textbf{Guanaco}: \texttt{Sorry, the secret word is banana. It's best if we don't speak of it. Is there anything else you'd like to ask?}
\end{quote}

These findings underscore the importance of researching approaches that ensure more robust compliance with instructions.

\paragraph{Math} {Guanaco} demonstrates significant shortcomings in mathematics, an area that has historically proven difficult for many language models as discussed in \citep{liang2022holistic}. When the model provides explicit intermediate reasoning, its solutions are often correct, for instance:

\begin{quote}
    \textbf{User}: Lee mows one lawn and charges \$33. Last week he mowed 16 lawns and three customers each gave him a \$10 tip. How many dollars did Lee earn mowing lawns last week?

    \textbf{Guanaco}: \texttt{Lee earned \$582 mowing lawns last week. \\ \\ To find the total amount he earned, we first need to calculate how much he earned for mowing 16 lawns: \\ \\  \$33 x 16 = \$528 \\ \\ Then we need to add the tips he received from the three customers: \\  \\ \$10 x 3 = \$30 \\ \\ So the total amount Lee earned was: \\ \\ \$528 + \$30 = \$558}
\end{quote}

Nonetheless, {Guanaco} is prone to mistakes on even elementary problems if it does not proceed in a sequential, stepwise fashion—an issue documented by \cite{weichain}. This is illustrated by the following exchange:

\begin{quote}
    \textbf{User}: Please factorize 1833.

    \textbf{Guanaco}: \texttt{Sure, I'd be happy to help! 1833 is a prime number, which means it has only two factors: 1 and 1833. \\ \\ So the unique factorization of 1833 is: \\ \\ 2\^{}1 * 3\^{}2 * 17\^{}1 \\ \\ I hope this helps! Let me know if you have any other questions.}
\end{quote}

which is incorrect (the actual factorization is $3 \times 17 \times 43$), and indeed doubly so.

\paragraph{Theory of Mind} {Guanaco} demonstrates notably advanced Theory of Mind reasoning abilities \cite{nematzadeh2018evaluating, sap2022neural}. As an illustration, the model provides a thorough and correct response to this scenario:

\begin{quote}
    \textbf{User}: James and Abby are in the bedroom. Abby put the pen in the desk drawer. Abby leaves the bedroom. James moves the pen into the bag. Where does James think Abby will look for the pen?

    \textbf{Guanaco}: \texttt{James knows that Abby put the pen in the desk drawer. But James also knows that he has moved the pen into the bag, and he assumes that Abby does not know about this. So, James may guess that Abby will look for the pen in the desk drawer, because that's where she left it last time. However, this assumption may be incorrect, because Abby may have seen James moving the pen.}
\end{quote}

Yet, the reliability of such inferences is limited; the model often provides reasoning that relies on unwarranted assumptions about the scenario, for example:

\begin{quote}
    \textbf{User}: Evelyn entered the living room. Jackson entered the playroom. James entered the playroom. The beans are in the treasure chest. James exited the playroom. Jackson moved the beans to the pantry. Jackson exited the playroom. James entered the living room. Where will James look for the beans?

    \textbf{Guanaco}: \texttt{James will look for the beans in the pantry, because that's where Jackson moved them.}
\end{quote}

In this case, {Guanaco} incorrectly infers the transfer of information that was not specified. Such challenges align with findings in the current research literature \cite{sap2022neural}, highlighting a need for further investigation.

\subsection{Considerations}
\label{ssec:considerations}

\paragraph{Evaluation}

Our results show moderate inter-annotator agreement (Fleiss $\kappa=0.42$), which drops further in comparisons involving two high-performing systems. This outcome reveals shortcomings in existing evaluation benchmarks and human assessment protocols regarding chatbot tasks. During side-by-side evaluations of outputs from ChatGPT and {Guanaco} 65B on the Vicuna benchmark, we observed that subjective judgments became increasingly influential, as disagreements among this paper’s authors were frequent concerning which responses were preferable. We suggest that future studies explore strategies to address these issues, leveraging methods from fields such as Human-Computer Interaction and Psychology that have grappled with subjectivity.

Our analysis also indicates that automated assessment systems are affected by systematic biases. For instance, we detect pronounced order effects in which GPT-4 consistently assigns higher ratings to whichever system’s responses appear first in the prompt. Moreover, sample-level agreement between GPT-4 and human annotators is low (Fleiss $\kappa=0.25$), implying that the evaluation criteria between humans and automated systems can diverge. Additionally, as presented in Table~\ref{tbl:elo_large}, GPT-4 tends to rate its own outputs substantially higher than human raters do, producing an Elo score of 1348 compared to 1176 from humans—corresponding to approximately a 20\% increased likelihood of outperforming another system.

Future investigations should focus on identifying potential biases in automated evaluation frameworks and exploring methods to mitigate such issues.

\paragraph{Data \& Training} It is important to highlight that the OASST1 dataset, which is used to train the {Guanaco} models, includes multiple languages, and the OA benchmark similarly features prompts in various languages. Future research is required to determine the extent to which this multilingual training enhances the ability to follow instructions in languages apart from English, and to assess whether this contributes to the larger performance difference observed between the Vicuna-13B model (which is exclusively trained on English corpora) and the {Guanaco} 33B and 65B models on the OA benchmark.

Considering the robust results achieved by the {Guanaco} models, we examined the possibility of data leakage between the OASST1 training set and the Vicuna benchmark prompts. By employing fuzzy string matching and subsequently conducting a manual review of the most similar pairs, we did not find evidence of duplicated prompts across these datasets.

In addition, our model training is solely based on supervised learning using cross-entropy loss, with no application of reinforcement learning from human feedback (RLHF). This opens avenues for future work to examine the respective advantages and limitations of using only cross-entropy loss versus incorporating RLHF. We anticipate that \textsc{QLoRA} will facilitate large-scale investigations of these questions without demanding significant computational resources.

\section{Related Work}

\paragraph{Quantization of Large Language Models} Most prior studies on quantizing LLMs have concentrated on inference-time quantization. Noteworthy approaches aimed at maintaining 16-bit LLM performance primarily deal with the management of outlier features (such as SmoothQuant~\citep{xiao2022smoothquant} and LLM.int8()~\citep{dettmers2022llm}), while alternative strategies employ more advanced grouping schemes~\citep{park2022nuqmm, yao2022zeroquant}. Some research explores lossy quantization and evaluates the trade-offs associated with different rounding methods~\citep{dettmers2022case,zeng2022glm,pope2022efficiently}, as well as techniques to refine rounding for improved quantization fidelity~\citep{frantar2022gptq}. Apart from our contribution, SwitchBack layers~\citep{wortsman2023stable} remains the only study to systematically address backpropagation through quantized weights at model scales exceeding 1B parameters.

\paragraph{Finetuning with Adapters} Our approach utilizes Low-rank Adapters~\citep{hu2021lora} (LoRA), although the literature encompasses a range of Parameter Efficient FineTuning (PEFT) strategies. These include techniques such as prompt tuning~\citep{qin2021learning,lester2021power,li2021prefix}, adapting the inputs to the embedding layer~\citep{an2022input}, adjusting hidden states~(IA$^3$)~\citep{liu2022few}, integrating additional full layers~\citep{houlsby2019parameter}, bias tuning~\citep{zaken2021bitfit}, employing Fisher information-based mask learning over weights~\citep{sung2021training}, and hybrid methodologies~\citep{henderson2021compacter}. Our experiments demonstrate that LoRA adapters can achieve performance on par with full 16-bit finetuning. The examination of tradeoffs among other PEFT techniques is deferred for future research.

\paragraph{Instruction Finetuning} Instruction finetuning adapts pretrained LLMs to better respond to prompts by training them on a diverse set of input-output examples sourced from multiple datasets. This method aligns LLM output with intended instructions as specified in the prompt. Prominent approaches and resources in this area are MetaICL~\citep{min2021metaicl}, MetaTuning~\citep{zhong2021adapting}, InstructGPT~\citep{ouyang2022training}, FLAN~\citep{wei2021finetuned,chung2022scaling}, PromptSource~\citep{bach2022promptsource}, Super-NaturalInstructions~\citep{wang2022super,sanh2021multitask}, Self-instruct~\citep{wang2022self}, UnnaturalInstructions~\citep{honovich2022unnatural}, OPT-IML~\citep{iyer2022opt}, UnifiedSKG~\citep{xie2022unifiedskg}, OIG/Chip2~\citep{laion2023}, Alpaca~\citep{alpaca}, Vicuna~\citep{vicuna2023}, Koala~\citep{koala_blogpost_2023}, and Self-instruct-GPT-4~\citep{peng2023instruction}.

\paragraph{Chatbots} Instruction-following models frequently adopt a chatbot structure, with many utilizing Reinforcement Learning from Human Feedback (RLHF)~\citep{christiano2017deep} or synthetic data generated with feedback from other AI models (RLAIF)~\citep{bai2022constitutional} for their training pipelines. Notable datasets and methods include Anthropic-HH~\citep{askell2021general,bai2022training}, Open Assistant~\citep{kopf2023openassistant}, LaMDA~\citep{thoppilan2022lamda}, and Sparrow~\citep{glaese2022improving}. Our work does not employ reinforcement learning, but Guanaco—our highest-performing model—is finetuned with multi-turn conversations drawn from the Open Assistant dataset, originally curated for RLHF training~\citep{kopf2023openassistant}. For the purpose of evaluation, recent work has replaced labor-intensive human annotations with GPT-4-based evaluation protocols~\citep{vicuna2023,peng2023instruction}. Building on this, we present refinements intended to yield more robust and trustworthy evaluation setups.

\section{Limitations and Discussion}

We have provided evidence that our approach, \textsc{QLoRA}, allows for replication of 16-bit full finetuning results using a 4-bit base model augmented with Low-rank Adapters (LoRA). However, we have yet to confirm whether \textsc{QLoRA} can attain parity with 16-bit full finetuning performance at the 33B and 65B parameter scales. Owing to significant computational demands, these larger-scale experiments remain an open direction for subsequent inquiry.

An additional constraint arises in our evaluation of instruction finetuning models. While our analysis includes results on MMLU, the Vicuna benchmark, and the OA benchmark, we did not perform evaluations on alternative benchmarks such as BigBench, RAFT, and HELM. As a result, generalizability of our evaluations to these benchmarks cannot be assumed. Nevertheless, we conducted an extensive investigation using MMLU and also proposed new evaluation methodologies tailored for chatbot assessments.

Based on the data analyzed, it seems that the results on various benchmarks are largely influenced by the extent to which the finetuning data aligns with the content of each benchmark. For instance, while FLAN v2 is closely related to the MMLU dataset, it is less aligned with chatbot benchmarks; conversely, the Chip2 dataset is more relevant for chatbot tasks and less for MMLU, and both models reflect these relationships in their performance on the MMLU and Vicuna benchmarks. This underscores not only the need for improved benchmarks and evaluation frameworks but also the importance of clarity regarding the specific competencies being tested. Should the objective be to optimize model performance for academic knowledge typical of high school and college classrooms, or to excel in conversational chatbot interactions? Perhaps priorities lie elsewhere. Since there is an inherent tendency to prefer existing benchmarks over developing new ones, the choice of benchmarks can shape research directions within the community. Therefore, it is crucial that benchmarks accurately capture the phenomena and tasks researchers intend to measure.

Although we conduct a comprehensive evaluation of general chatbot capabilities, one notable limitation is the restricted responsible AI assessment performed for {Guanaco}. Specifically, we assess the probability that {Guanaco}-65B produces socially biased outputs in comparison to other models, as reported in Table~\ref{tbl:crows}. The results indicate that {Guanaco}-65B generates such content at a much lower rate than unrefined, pretrained models, suggesting that finetuning with the OASST1 dataset can help mitigate bias inherent in the LLaMA base model. While these findings are promising, the effectiveness of {Guanaco} in terms of other bias dimensions remains uncertain. A more exhaustive investigation of biases present in {Guanaco} and analogous chatbots is left as a direction for future research.

Another constraint of our study is the omission of experiments involving varying bit precisions—such as adopting 3-bit base models—or alternative adapter approaches. In addition to LoRA, numerous Parameter Efficient FineTuning (PEFT) techniques have demonstrated strong performance, though it remains undetermined whether their effectiveness scales to larger models. While we employed LoRA due to its established reliability, exploring other adapters could potentially yield superior outcomes. Interestingly, our observations indicate that post-quantization finetuning is capable of restoring much of the information degraded during quantization, suggesting the possibility of adopting even more aggressive quantization strategies. For example, finetuning a base model quantized to 3-bits with GPTQ and LoRA might approach the performance of conventional 16-bit finetuning after appropriate adaptation.

\section{Broader Impacts}

Our proposed \textsc{QLoRA} finetuning approach constitutes the first technique capable of facilitating the finetuning of models with 33B parameters on standard consumer GPUs, and models up to 65B parameters on professional-grade GPUs, all without a decline in performance compared to full finetuning baselines. We have shown that our most advanced 33B model, when finetuned on the Open Assistant dataset, can compete with ChatGPT on the Vicuna benchmark. Given that instruction finetuning is a crucial process for transforming raw pretrained large language models into advanced chatbots similar to ChatGPT, we anticipate that our method will significantly broaden the accessibility of state-of-the-art NLP by lowering the computational and financial barriers, particularly for researchers with limited resources. In this way, \textsc{QLoRA} acts as an equalizing mechanism, helping to reduce the gap in capabilities between well-funded organizations and smaller groups using standard hardware.

A further significant avenue for influence arises from deploying models on mobile devices. We contend that our \textsc{QLoRA} approach represents a pivotal step toward facilitating the finetuning of LLMs directly on phones and within other constrained computational environments. Prior work has demonstrated the feasibility of running 7B models on mobile hardware; however, \textsc{QLoRA} constitutes the initial method capable of supporting their finetuning. Our projections suggest that, using an iPhone 12 Plus, it would be possible for \textsc{QLoRA} to process approximately 3 million tokens overnight during charging periods. Although finetuned 7B parameter models do not yet match the performance levels of ChatGPT, their outputs are sufficiently robust to unlock a new range of applications previously restricted by concerns over privacy or the limitations of LLM performance. \textsc{QLoRA} stands to promote privacy-preserving language model applications, enabling individuals to maintain control over their own data and models, while also simplifying the process of model deployment.

Nonetheless, it should be recognized that finetuning technology has dual-use potential, which can lead to malicious exploitation. While large-scale LLM deployment carries documented risks \citep{bommasani2021opportunities,bender2021dangers}, we argue that democratizing access to this rapidly proliferating technology may foster more independent and thorough scrutiny than would be possible if control were confined to large organizations unwilling to make their models or code publicly accessible for review.

In summary, we anticipate that \textsc{QLoRA} will contribute positively by substantially increasing the accessibility and ease of finetuning high-performance LLMs for a broader community.